\documentclass[11pt,a4paper]{article}

% ========================
% ESSENTIAL PACKAGES
% ========================

% Language and encoding
\usepackage[utf8]{inputenc}
\usepackage[T1]{fontenc}
\usepackage[english]{babel}
%\usepackage{natbib}

% Math packages
\usepackage{amsmath}        % Enhanced math environments
\usepackage{amssymb}        % Additional math symbols
\usepackage{amsthm}         % Theorem environments
\usepackage{amsfonts}       % Additional math fonts
\usepackage{mathtools}      % Extension to amsmath
\usepackage{bm}             % Bold math symbols
\usepackage{dsfont}         % Double-struck fonts (for sets like ℝ, ℕ)
\usepackage{mathrsfs}       % Script fonts for math

% Graphics and figures
\usepackage{graphicx}       % Include graphics
\usepackage{float}          % Better float placement
\usepackage{subcaption}     % Subfigures
\usepackage{tikz}           % Drawing diagrams
\usepackage{pgfplots}       % Plotting
\pgfplotsset{compat=1.18}

% Page layout and formatting
\usepackage[margin=1in]{geometry}
\usepackage{fancyhdr}       % Custom headers and footers
\usepackage{titlesec}       % Customize section titles
\usepackage{enumitem}       % Customize lists
\usepackage{parskip}        % Paragraph spacing

% Colors and highlighting
\usepackage{xcolor}
\usepackage{mdframed}       % Framed environments
\usepackage{tcolorbox}      % Colored boxes

% Tables
\usepackage{array}          % Extended array environments
\usepackage{booktabs}       % Professional tables
\usepackage{multirow}       % Multi-row cells

% References and links
\usepackage{hyperref}       % Hyperlinks
\usepackage{cleveref}       % Smart cross-references

% Additional useful packages
\usepackage{cancel}         % Cancel terms in equations
\usepackage{siunitx}        % SI units
\usepackage{listings}       % Code listings (if needed)

% ========================
% CUSTOM COMMANDS
% ========================

% Common math sets
\newcommand{\R}{\mathbb{R}}
\newcommand{\N}{\mathbb{N}}
\newcommand{\Z}{\mathbb{Z}}
\newcommand{\Q}{\mathbb{Q}}
\newcommand{\C}{\mathbb{C}}

% Common operators
\DeclareMathOperator{\lcm}{lcm}
\DeclareMathOperator{\Tr}{Tr}
\DeclareMathOperator{\rank}{rank}
%\DeclareMathOperator{\span}{span}
%\DeclareMathOperator{\null}{null}

% Shortcuts for common expressions
\newcommand{\abs}[1]{\left|#1\right|}
\newcommand{\norm}[1]{\left\|#1\right\|}
\newcommand{\inner}[2]{\left\langle #1, #2 \right\rangle}
\newcommand{\ceil}[1]{\left\lceil #1 \right\rceil}
\newcommand{\floor}[1]{\left\lfloor #1 \right\rfloor}

% Derivatives and integrals
\newcommand{\dd}[1]{\,\mathrm{d}#1}
\newcommand{\dv}[2]{\frac{\mathrm{d}#1}{\mathrm{d}#2}}
\newcommand{\pdv}[2]{\frac{\partial #1}{\partial #2}}

\renewcommand{\hat}{\widehat}
\renewcommand{\tilde}{\widetilde}

% ========================
% THEOREM ENVIRONMENTS
% ========================

\theoremstyle{definition}
\newtheorem{definition}{Definition}[section]
\newtheorem{example}{Example}[section]
\newtheorem{exercise}{Exercise}[section]

\theoremstyle{plain}
\newtheorem{theorem}{Theorem}[section]
\newtheorem{lemma}{Lemma}[section]
\newtheorem{corollary}{Corollary}[section]
\newtheorem{proposition}{Proposition}[section]
\newtheorem{assumption}{Assumptio}[section]

\theoremstyle{remark}
\newtheorem{remark}{Remark}[section]
\newtheorem{note}{Note}[section]

\newcommand{\supp}{\mathrm{supp}}
\newcommand{\op}{\mathrm{op}}

\newcommand{\mX}{\mathcal X}
\newcommand{\mD}{\mathcal D}
\newcommand{\bP}{\mathbb P}
\newcommand{\lp}{L^2(\bP)}
\newcommand{\lx}{L^2(\mX)}
\newcommand{\ltwo}{L^2(\mX)}
\newcommand{\ud}{\mathrm d}
\newcommand{\mH}{\mathcal H}
\newcommand{\mA}{\mathcal A}

% ========================
% COLORED BOXES
% ========================

% Important theorem box
\newtcolorbox{importantbox}{
	colback=blue!5!white,
	colframe=blue!75!black,
	title=Important
}

% Warning/Note box
\newtcolorbox{notebox}{
	colback=yellow!10!white,
	colframe=orange!75!black,
	title=Note
}

% Example box
\newtcolorbox{examplebox}{
	colback=green!5!white,
	colframe=green!75!black,
	title=Example
}

% ========================
% DOCUMENT SETUP
% ========================

% Page style
\pagestyle{fancy}
\fancyhf{}
\rhead{\thepage}
\lhead{\leftmark}

% Title formatting
\title{Functional Bandit Notes}
\author{Your Name}
\date{\today}

% Hyperlink setup
\hypersetup{
	colorlinks=true,
	linkcolor=blue,
	filecolor=magenta,      
	urlcolor=cyan,
	citecolor=red
}

\begin{document}

The main purpose of this document is to give a technical summary of existing approaches towards (offline) functional regression,
and remark on their strengths/weakness. When appropriate, I will also remark on possibilities to bandit problems (mostly confidence intervals).

Let $\mX$ be a compact Euclidean set and $\mD=L^2(\mX)$ be the set of all square integrable functions on $\mX$.
Let $\mathbb P$ be the distribution on $\mD$. The model is 
$$
y_i = \langle f_i, \beta_0\rangle + \varepsilon_i
$$
where $\{\varepsilon_i\}_{i=1}^n$ are centered, exogenous sub-Gaussian random variables, $\beta_0\in\mathcal F\subseteq \mD$ is an unknown regression model and $f_1,\cdots,f_N\overset{i.i.d.}{\sim}\mathbb P$.
The inner product is taken with respect to the Lesbesgue measure on $\mX$:
$$
\langle f_i, \beta_0\rangle := \int_{\mX} f_i(u)\beta_0(u)\ud u.
$$
Sometimes, we will place additional assumptions on $\mathcal F$, or restrict the support of $\bP$.
Errors are measured in 
$$
\|\hat \beta_n-\beta_0\|_{\lp}^2 := \int_{\mD} \langle f, \hat \beta_n-\beta_0\rangle^2\ud\bP(f).
$$

Let $K$ be the RKHS and recall the notation that $\tilde f = L_{K^{1/2}}f$, $\tilde\beta_0=L_{K^{-1/2}}\beta_0$, such that $\|\tilde\beta_0\|_2 < +\infty$.
Let $T=\mathbb E[\tilde f\otimes \tilde f]$ be the population covariance operator, which is decomposed as $T=\sum_{j=1}^{\infty}s_j\psi_j\otimes\psi_j$ for an orthonormal basis $\{\psi_j\}_{j=1}^{\infty}$ of $L_2(\mX)$,
and $s_j\asymp j^{-2r}$.

\paragraph{The self-similarity condition.}
For a test function $f$ and $\tilde f=L_{K^{1/2}}f$, it satisfies the self-similarity condition with respect to $\tilde\beta_0$, if there exists $c>0$ such that for all $1\leq k\leq \sqrt{n}$, 
$$
\frac{1}{c}\leq \left|\langle \tilde\beta_0, (k T+k^{-1}I)^{-1}\tilde f\rangle\right| \leq C.
$$

\begin{remark}
We only need the inequality to hold over a geometric grid $k=2^{\ell}$ for $\ell=0,1,\cdots,\log_2(\sqrt{n})$.
\end{remark}

\begin{remark}
Unlike the anisotropic approach, this framework does \emph{not} require self-similarity to be defined on the sample covariance operator, because the concentration from $T_n$ to $T$ is supposedly higher-order. In this anisotropic approach because we can only do anisotropic regularization on $\{\hat\psi\}$ this becomes an issue. In Lepski, for a particular $\tilde f$, regularization is still isotropic (albeit different $\lambda$ is chosen for different $\tilde f$), so the population covariance operator is sufficient.
\end{remark}

\begin{remark}
The scaling in the condition is non-homogeneous ($k$ versus $k^2$) so it might look strange. This condition, however, stems directly from the \emph{exact} bias expression
$$
\langle\tilde f, \beta_\lambda-\tilde\beta_0\rangle = \sqrt{\lambda}\cdot \langle \tilde\beta_0,\sqrt{\lambda}(T+\lambda)^{-1}\tilde f\rangle.
$$
If we want the bias to be upper and lower bounded by $\sqrt{\lambda}$, then the inner product must stay at constant level. If $\lambda=1/k^2$ for some $1\leq k\leq \sqrt{n}$, then
\begin{align*}
\sqrt{\lambda}(T+\lambda)^{-1}\tilde f &= (T+k^{-2})^{-1} k^{-1} \tilde f = (kT+k^{-1}I)^{-1} \tilde f.
\end{align*}
\end{remark}

\bibliographystyle{plain}
\bibliography{refs}

\end{document}